\section{Threats to Validity}
\label{sec:threats}

\textbf{Construct Validity}:
The primary threat to construct validity relates to whether APFD accurately reflects real-world testing efficiency.
While APFD is the standard evaluation metric in regression testing and competition benchmarks, ranking ties can artificially distort APFD values.
We mitigate this by conducting $30$ independent random trials with randomized tie-breaking and sub-trial sampling, reporting both the mean and standard deviation.
Additionally, we report ROC-AUC to evaluate threshold-free classification capability.

\textbf{Internal Validity}:
Internal validity concerns potential biases introduced during data processing and model optimization.
All comparative models were trained under identical preprocessing pipelines and consistent random seeds ($42$).
To avoid data leakage, per-channel normalization statistics and SWA checkpoints were calculated strictly on the training partition.

\textbf{External Validity}:
A potential threat to external validity is whether our findings generalize beyond the BeamNG.tech simulator and the SensoDat benchmark.
We mitigated this threat by evaluating SE2RoadNet across eight diverse benchmarks (including SDC-Scissor, OOB-0-1, OOB-0-3, and SDC-Travel) spanning different road generators and driving agents.
While our formulation is grounded in general Newtonian vehicle dynamics, physical validation on physical autonomous vehicles in real-world environments remains a subject for future investigation.
